\section{Muestra y muestreo}

La muestra corresponde a un subconjunto de la población de estudio y consiste en el número de casos que realmente serán estudiados. Esta debe ser representativa para garantizar la validez de los resultados. Por otro lado, el muestreo se refiere al proceso aplicado por el investigador para seleccionar a los participantes dentro de la población \cite{muestra}.

Existe el muestreo probabilístico (donde cada miembro de la población tiene una probabilidad de ser escogido) y no probabilístico (donde la selección se basa en la conveniencia, el juicio del investigador u otros criterios) \cite{muestra}. 

Para esta investigación se realiza un muestreo no probabilístico intencional, donde los elementos seleccionados para la muestra son elegidos por el criterio del investigador. Esto se justifica, ya que, al existir una infinitud de autómatas celulares dadas las restricciones en la población, es necesario delimitar y tratar de replicar los comportamientos. Además, dado que cada autómata celular resuelve un problema, es necesario identificar de forma precisa el problema que se está resolviendo \cite{computationalBiology2016}.